\begin{table}[H]
\centering
\caption{Crecimiento departamental de remuneraci\'on y PIB por trabajador, 2010--2024pr}
\label{tab:pib_geih_productividad_departamento_remuneracion_crecimientos}
\scriptsize
\begin{tabular}{lrrrrr}
\toprule
Departamento & Crec. rem./trab. & Puesto & Crec. PIB/trab. & Puesto & Dif. puestos \\
\midrule
Risaralda & 2,83\% & 1 & 3,14\% & 3 & -2 \\
Cundinamarca & 2,39\% & 2 & -0,32\% & 23 & -21 \\
Caldas & 2,30\% & 3 & 2,14\% & 6 & -3 \\
Bogotá D.C. & 1,90\% & 4 & 1,98\% & 8 & -4 \\
Boyacá & 1,66\% & 5 & 1,16\% & 15 & -10 \\
Antioquia & 1,65\% & 6 & 1,27\% & 13 & -7 \\
Tolima & 1,55\% & 7 & 2,04\% & 7 & +0 \\
Quindío & 1,50\% & 8 & 3,78\% & 1 & +7 \\
Valle del Cauca & 1,47\% & 9 & 2,15\% & 5 & +4 \\
Meta & 1,43\% & 10 & 0,05\% & 20 & -10 \\
Bolívar & 1,20\% & 11 & 1,69\% & 10 & +1 \\
Cauca & 1,20\% & 12 & 1,44\% & 12 & +0 \\
Atlántico & 1,09\% & 13 & 0,46\% & 18 & -5 \\
Chocó & 1,09\% & 14 & 2,95\% & 4 & +10 \\
Norte de Santander & 1,06\% & 15 & 1,23\% & 14 & +1 \\
Caquetá & 0,97\% & 16 & 3,71\% & 2 & +14 \\
Huila & 0,95\% & 17 & -0,13\% & 22 & -5 \\
Córdoba & 0,79\% & 18 & 0,90\% & 17 & +1 \\
Nariño & 0,63\% & 19 & 0,97\% & 16 & +3 \\
Santander & 0,62\% & 20 & 1,71\% & 9 & +11 \\
Sucre & 0,35\% & 21 & 1,67\% & 11 & +10 \\
Magdalena & 0,24\% & 22 & 0,16\% & 19 & +3 \\
Cesar & 0,16\% & 23 & -0,03\% & 21 & +2 \\
La Guajira & -1,84\% & 24 & -0,80\% & 24 & +0 \\
\bottomrule
\end{tabular}
\caption*{\footnotesize Nota: tasas de crecimiento anualizadas. La diferencia de puestos corresponde al puesto en crecimiento de la remuneraci\'on por trabajador menos el puesto en crecimiento del PIB por trabajador; un valor positivo indica que el departamento ocupa una posici\'on m\'as baja en crecimiento de remuneraci\'on que en crecimiento de productividad. Fuente: c\'alculos propios con DANE y GEIH.}
\end{table}